% Minimal preamble for P2 cards taken from sources/comparability-partitions-2408.13523v2/source0.tex
% (arXiv:2408.13523v2, "Partitioning perfect graphs into comparability graphs").
% Every line below is copied mechanically from that file; nothing was invented, renamed or redefined.
%   line 2    : \documentclass[12pt]{article}
%   line 3    : \usepackage{amssymb,latexsym}
%   line 4    : \usepackage{amsmath}
%   line 5    : \usepackage[dvipdfm]{graphicx}  -- the dvipdfm driver option is dropped, because the
%               assembler forces the xelatex engine and a forced driver option contradicts it.
%   line 16   : \usepackage{amsfonts}
%   line 17   : \usepackage{amsthm}
%   lines 20-30: the author's theorem declarations
%   lines 32-41: the author's own macro definitions
% Omitted source preamble lines and why they are not needed for these body spans:
%   line 1 (provenance comment), lines 6-11 (tikz and its libraries: no drawing in the retained spans),
%   lines 7 and 9 (verbatim), lines 12-13 (page geometry), line 14 (rawfonts),
%   line 15 (psfrag), line 18 (hyperref with the pagebackref option: the retained spans print no
%   citation anchors and the assembly adds its own hyperref machinery), lines 43-61 (the author's
%   backref neutralisation and his margin-note macros, which depend on the omitted hyperref).
% The retained spans of the body use only \em, \bf, \cal, amsthm environments, amsmath displays and
% the macros defined in lines 32-41 below.
\documentclass[12pt]{article}
\usepackage{amssymb,latexsym}
\usepackage{amsmath}
\usepackage{graphicx}
\usepackage{amsfonts}
\usepackage{amsthm}

\newtheorem{theorem}{Theorem}%[section]
\newtheorem{problem}[theorem]{Problem}%[section]
\newtheorem{observation}[theorem]{Observation}
\newtheorem{lemma}[theorem]{Lemma}
\newtheorem{corollary}[theorem]{Corollary}
\newtheorem{definition}{Definition}
\newtheorem{example}[theorem]{Example}
\newtheorem{conjecture}{Conjecture}%[section]
\newtheorem{proposition}[theorem]{Proposition}
\newtheorem{claim}[theorem]{Claim}
\newtheorem{remark}[theorem]{Remark}

\def\qed{\ifhmode\unskip\nobreak\fi\quad\ifmmode\Box\else$\Box$\fi}
\def\pont{\hspace{-6pt}{\bf.\ }}
\def\qed{\hfill $\Box$}

\def\R{{\mbox{\boldmath$\mathcal R$}}}
% the author's second definition of \R, left commented out exactly as in the source
\def\F{{\mathcal F}}
\def\P{{\mathcal P}}

\def\newpage{\vfill\eject}
\title{A squared-loglog upper bound for the comparability partition number of the interval graph G\_n}
\author{András Gyárfás, Márton Marits, Géza Tóth}
\date{arxiv:2408.13523v2; CC BY 4.0}
\begin{document}
\maketitle
\noindent\textit{Editorial note.} The counted unit is the article's own Theorem ub at source lines 468-470 together with the author's own proof environment at source lines 609-618, whose header is Proof of Theorem ub; the proof is not adjacent to the statement, because the author places the preparatory definition of F\_n at source lines 475-479, the recurrence Lemma rec with its long proof at source lines 484-598 and the corollary p' at source lines 602-605 in between, and the pairing therefore rests on the author's own named header rather than on a later-first-proof rule.\par\medskip
\noindent\textit{Editorial note.} The retained context at source lines 128-129, 391-395, 409-411, 475-479 and 602-605 is the article's own definitions of the parameters p(G) and c(G), of the interval graph G\_n and of the numbers c(n), p(n), F\_n and p'(n), together with two of its own corollaries: the bound for the partition number of the shift graph S\_n and the quadratic bound for p'(n) at powers of two. Both corollaries are retained as quoted prerequisites of the counted proof, not as counted units.\par\medskip
\noindent\textit{Editorial note.} The retained context at source lines 416-419 and 602-605 carries no citation and no cross-reference. The body of Lemma rec at source lines 484-598, which is the interval that is publicly occupied for this source, is deliberately not retained.\par\medskip
\noindent\textit{Editorial note.} The source preamble cannot be retained as it stands: it loads rawfonts, psfrag, verbatim, tikz with several libraries, hyperref with the pagebackref option and a dvi driver option for graphicx, none of which the retained body spans use, and it redefines backref commands that only exist because of that hyperref option. The retained preamble file therefore carries, line for line, source line 2, source line 3, source line 4, source line 16, source line 17, the theorem declarations of source lines 20-30 and the author macro definitions of source lines 32-41; the only textual change anywhere is the removal of the dvipdfm driver option from the graphicx line, because the assembly forces the xelatex engine. No new macro, package or declaration was introduced.\par\medskip
\noindent\textit{Editorial note.} The article's theorems co-triangulated, unit, gensplit, bipline and loglog, its lemma rec, its corollaries doblesh and shift as counted statements, and its proposition H9 are not counted here.\par\medskip
\noindent\textit{Allowed standard prerequisites (editorial).}\begin{itemize}
\item The partition number p(G) and the cover number c(G). p(G) is the minimum m such that the edge set of G is the union of m edge-disjoint comparability subgraphs of G, and c(G) is the same number without the edge-disjointness requirement; retained at source lines 128-129.. The counted theorem is an upper bound for p(n)=p(G\_n), so the parameter and its edge-disjointness requirement must be fixed before the statement can be read.
\item The interval graph G\_n and the number p(n). For n at least 2, G\_n is the intersection graph of the n choose 2 closed intervals whose endpoints are among n consecutive positive integers, and the article abbreviates c(n)=c(G\_n) and p(n)=p(G\_n); retained at source lines 391-395.. Identifies the graph whose partition number the counted theorem bounds, and shows that p(n) is non-decreasing in n, which the proof uses to reduce to the powers of two.
\item The shift graph S\_n and the bound on its partition number. S\_n is the subgraph of the touching edges of G\_n, its chromatic number equals the ceiling of the logarithm of n, and the article records at source lines 416-419 the corollary that p(S\_n) is at most the ceiling of the logarithm of the ceiling of the logarithm of n, with a partition into bipartite graphs; retained at source lines 409-411 and 416-419.. The proof splits the edge set of G\_n into the touching part S\_n and the rest F\_n and quotes exactly this bound for the first part.
\item The graph F\_n and the number p'(n). F\_n is obtained from G\_n by deleting all touching edges without deleting their endpoints, and p'(n) is defined as p(F\_n); retained at source lines 475-479.. The proof uses that the edge set of G\_n is the disjoint union of the edge sets of S\_n and F\_n, which is what allows the two partitions to be combined.
\item The corollary bounding p'(n) at the powers of two. For every n of the form 2\textasciicircum{}\{2\textasciicircum{}k\} with k at least 1 the article proves the chain of inequalities p'(n)=p'(2\textasciicircum{}\{2\textasciicircum{}\{k-1\}\}) <= p(2\textasciicircum{}\{2\textasciicircum{}\{k-1\}\})+k-1+2 <= k choose 2 + 2k - 1; retained at source lines 602-605.. The proof substitutes exactly this quadratic bound for p'(n) at n=2\textasciicircum{}\{2\textasciicircum{}k\} to obtain the stated estimate.
\end{itemize}\medskip
Let $p(G)$ denote the minimum number $m$ such that $E(G)=\cup_{i=1}^m E(G_i)$ where the $G_i$-s are edge-disjoint comparability subgraphs of $G$.
Dropping  the condition that the $G_i$-s are edge-disjoint, we define $c(G)$ as the smallest number of comparability subgraphs of $G$ needed to {\em cover} $E(G)$.

For $n\ge 2$, let $G_n$ be the interval graph defined as the intersection graph
of the $n\choose 2$ closed intervals with endpoints among $n$ consecutive positive integers,
called {\em base points}. Note that $|V(G_n)|={n\choose 2}$.
For convenience, set $c(n)=c(G_n),p(n)=p(G_n)$.
Since $G_2$ has one vertex and no edge, we set $c(2)=p(2)=0$.

Let $S_n$ denote the subgraph of touching edges of $G_n$.
The graph $S_n$ is often referred as the {\em shift graph}.
It is well-known that $\chi(S_n)=\lceil\log(n)\rceil$

\begin{corollary}\label{shift}
  $p(S_n)\le \lceil \log(\chi(S_n))\rceil = \lceil\log(\lceil\log(n)\rceil)\rceil$
  and the partition can be done by bipartite graphs.
\end{corollary}

Before the proof, we need some preparations.
Let $F_n$ be the graph obtained from $G_n$ by removing all touching edges of $G_n$,
i.e. edges of $S_n$ (without removing their endpoints).
Let
$p'(n)=p(F_n)$.

\begin{corollary}\label{p'}
  For any  $n= 2^{2^k}$ with $k\ge 1$ we have
  $$p'(n)=p'(2^{2^{k-1}})\le p(2^{2^{k-1}})+k-1+2 \le \dots \le {k\choose 2}+2k-1.$$
\end{corollary}

\begin{theorem}\label{ub} For any $n\ge 4$ we have
  $$p(n)\le {1\over 2}(\lceil\log(\lceil\log(n)\rceil)\rceil)^2+{5\over 2}\lceil\log(\lceil\log(n)\rceil)\rceil-1.$$
\end{theorem}

\begin{proof}[Proof of Theorem \ref{ub}]
It is enough to prove the statement for $n=2^{2^k}$ where $k\ge 2$.
Since $p(n)$ is non-decreasing, this implies the statement for $2^{2^{k-1}}<n < 2^{2^k}$ as well.

Let $n=2^{2^k}$. Observe that the edge set of $G_n$ is the disiont union of the edge set of $S_n$ and $F_n$.
Combining the partitions given by Corollary \ref{shift} and \ref{p'} respectively we get that

$$p(n) \le p'(n)+ p(S_n)\le {k\choose 2}+2k-1+ \lceil\log(\lceil\log(n)\rceil)\rceil={k\choose 2}+3k-1$$
$$\le {1\over 2}(\lceil\log(\lceil\log(n)\rceil)\rceil)^2+{5\over 2}\lceil\log(\lceil\log(n)\rceil)\rceil-1.$$
finishing the proof of Theorem \ref{ub}.
\end{proof}

\end{document}
